\documentclass[10pt,a4paper]{article}

% ---------- Packages ----------
\usepackage[T1]{fontenc}
\usepackage[utf8]{inputenc}
\usepackage{lmodern}
\usepackage[french]{babel}
\usepackage[margin=1.4cm, top=0.7cm, bottom=0.7cm]{geometry}
\usepackage{graphicx}
\usepackage{enumitem}
\usepackage{titlesec}
\usepackage[hidelinks]{hyperref}

\pagestyle{empty}
\linespread{0.95}
\setlength{\parindent}{0pt}
\setlength{\parskip}{0pt}
\frenchbsetup{StandardLists=true}

% ---------- Style des sections (noir et blanc) ----------
\titleformat{\section}{\large\bfseries\scshape}{}{0em}{}[\vspace{-0.55em}\rule{\linewidth}{0.5pt}]
\titlespacing*{\section}{0pt}{6pt}{3pt}

% ---------- Listes ----------
\setlist[itemize]{leftmargin=1.2em, itemsep=0.5pt, parsep=0pt, topsep=1pt, label=\textbullet}

% ---------- Commandes ----------
% \entry{Titre}{Date}{Sous-titre / lieu}
\newcommand{\entry}[3]{%
  \textbf{#1} \hfill \textbf{#2}\\
  \textit{#3}\par\vspace{-2pt}
}
% \entryone{Titre}{Date}  (sans sous-titre)
\newcommand{\entryone}[2]{%
  \textbf{#1} \hfill \textbf{#2}\par\vspace{-2pt}
}

\begin{document}

% ---------- En-tête avec photo ----------
% Place ta photo dans le même dossier sous le nom photo.jpg (ou change le nom ci-dessous)
\newcommand{\photo}{%
  \IfFileExists{photo.jpg}{%
    \includegraphics[width=2.5cm,height=3.1cm]{photo.jpg}}{%
    \fbox{\parbox[c][3.1cm][c]{2.5cm}{\centering\small Photo}}}%
}
\noindent
\begin{minipage}[c]{0.77\textwidth}
  {\huge\bfseries Laora Aimi}\\[5pt]
  {\Large\bfseries Stage Ingénieure Support Client -- Programmes}\\[1pt]
  {\large Stage de 6 mois, à partir de mars 2027}\\[6pt]
  \href{mailto:laora.aimi@gmail.com}{laora.aimi@gmail.com} \ $|$ \ +33 7 81 74 44 45 \ $|$ \ Laval, France\\
  \href{https://www.linkedin.com/in/laora-aimi-493285253}{linkedin.com/in/laora-aimi-493285253} \ $|$ \
  \href{https://github.com/Laoraa}{github.com/Laoraa}
\end{minipage}%
\hfill
\begin{minipage}[c]{0.2\textwidth}
  \raggedleft\photo
\end{minipage}

\vspace{-2pt}
% ---------- Profil ----------
\section{Profil}
Étudiante en Master Motorsport Engineering Performance à l'ESTACA, formée à l'acquisition et à l'analyse de données de piste (Motec), avec une solide formation en informatique et data : bases de données, pipelines automatisés, outils de visualisation. Passionnée de sport automobile et très mobile, je suis motivée à me déplacer sur les circuits, en France comme à l'étranger, pour assister les équipes en piste.

% ---------- Formation ----------
\section{Formation}

\entryone{ESTACA -- Master Motorsport Engineering Performance, Laval}{2026 -- 2027}
\begin{itemize}
  \item \textbf{Données de piste :} acquisition par capteurs, analyse de performance (Motec), analyse sur simulateur (WINTAX), traitement de données sous MATLAB
  \item \textbf{Opérations en piste :} stratégie de course (Trackside), analyse du comportement pilote, pilotage d'une Formule 4 (FFSA)
\end{itemize}

\vspace{2pt}
\entryone{ENSEIRB-MATMECA -- Diplôme d'ingénieur en informatique, Bordeaux}{2023 -- 2026}
\begin{itemize}
  \item Spécialisation en génie logiciel et analyse de données
  \item Semestre d'échange à \textbf{Keio University}, Tokyo (avr. -- août 2025) : intelligence artificielle et vision par ordinateur
\end{itemize}

% ---------- Expérience professionnelle ----------
\section{Expérience professionnelle}

\entryone{BeTomorrow -- Stage de 6 mois, Bordeaux}{2026}
\begin{itemize}
  \item Développement d'un système de rapprochement de produits par IA pour E.Leclerc.
  \item Infrastructure cloud serverless sur GCP (Docker, Cloud Run Jobs) offrant au client une autonomie complète.
  \item Conception et déploiement de tableaux de bord Data Studio interactifs pour le reporting client.
\end{itemize}

\vspace{2pt}
\entryone{INRIA -- Stage de 3 mois, Bordeaux}{2025}
\begin{itemize}
  \item Construction de modèles prédictifs pour optimiser les performances d'applications en anticipant les opérations d'E/S.
  \item Amélioration de la précision des prédictions et de l'efficacité des applications.
\end{itemize}

\vspace{2pt}
\entryone{LORIA -- Stage de 6 semaines, Nancy}{Juin 2023}
\begin{itemize}
  \item Développement d'une application Android générant des images sphériques à partir de photos de smartphone.
\end{itemize}

% ---------- Projets sport auto & data ----------
\section{Projets sport automobile \& data}

\entryone{Étude de cas data : comparaison intra-équipe McLaren (GP de Hongrie 2025)}{}
\begin{itemize}
  \item Analyse Norris vs Piastri (temps par secteur, styles de pilotage, dégradation des pneus), avec modélisation des relais et visualisations stratégiques dans le style d'un débriefing d'ingénieurs.
\end{itemize}

\vspace{2pt}
\entryone{Rapport d'analyse de la saison de F1 2024}{\href{https://github.com/Laoraa/F1_Data_Analysis}{\normalfont\small github.com/Laoraa/F1\_Data\_Analysis}}
\begin{itemize}
  \item Rapport \LaTeX{} complet sur la saison 2024 (circuits, stratégies de course, tendances de performance)
\end{itemize}

% ---------- Projets académiques ----------
\section{Projets académiques}

\entryone{Convertisseur audio vers partition (Python)}{Équipe de 5 -- 2 mois}
\begin{itemize}
  \item Développement d'un système d'IA capable de convertir des enregistrements audio en fichiers MIDI et en partitions musicales.
\end{itemize}

\vspace{2pt}
\entryone{Voitures autonomes (Python)}{Équipe de 4 -- 4 mois}
\begin{itemize}
  \item Cheffe de projet : réalisation d'un programme permettant aux voitures de cartographier leur environnement.
\end{itemize}

% ---------- Compétences ----------
\section{Compétences \& langues}
\begin{itemize}[label={}, leftmargin=0pt, itemsep=2pt]
  \item \textbf{Programmation :} Python, C, C++, SQL, Java, TypeScript, HTML/CSS, MATLAB
  \item \textbf{Data \& analyse :} pandas, NumPy, scikit-learn, matplotlib, seaborn, modélisation prédictive, machine learning (régression, classification, clustering), visualisation de données
  \item \textbf{Outils :} Motec, WINTAX, Git/GitHub, Docker, Terraform, GCP (Cloud Run, GCS), Data Studio / Power BI
  \item \textbf{Langues :} français (langue maternelle), anglais (C1, TOEIC 950/990), allemand (A2), japonais (débutant)
\end{itemize}

% ---------- Centres d'intérêt ----------
\section{Centres d'intérêt}
Escrime (niveau national, arbitre régional) \ $|$ \ Esport : manager et data analyst d'une leading team sur \textit{Pok\'emon UNITE} \ $|$ \ Crochet, dessin, conception lumière

\end{document}
